\documentclass{article}
\usepackage[utf8]{inputenc}
\usepackage{amsmath}
\usepackage{geometry}
\geometry{a4paper, margin=1in}

\title{Mathematical Approaches to Superluminal Phenomenon \\ \large Simulating Astrophysical Shockwaves \& Jets}
\author{Solar System Simulation Research}
\date{February 24, 2026}

\begin{document}

\maketitle

This document outlines the primary computational mathematics frameworks used in 2024-2025 to simulate and approximate superluminal phenomena in astrophysical environments.

\section*{1. Apparent Superluminal Motion (Kinematics)}
The foundational analytical approach to modeling ``faster-than-light'' observation is a geometric projection effect driven by finite light-travel time.

If a plasma blob in an AGN jet moves at true velocity $v = \beta c$ at an angle $\theta$ relative to the observer's line of sight, the apparent transverse velocity $v_{app} = \beta_{app} c$ is:
\begin{equation}
    \beta_{app} = \frac{\beta \sin\theta}{1 - \beta \cos\theta}
\end{equation}

\section*{2. General Relativistic Magnetohydrodynamics (GRMHD)}
Used to model the macroscopic fluid dynamics of the jet launching and accretion disk surrounding the SMBH in curved spacetime.

GRMHD solves the conservation of particle number, energy, and momentum, coupled with Maxwell's equations in the ideal MHD limit:
\begin{align}
    \nabla_\mu (\rho u^\mu) &= 0 \\
    \nabla_\mu T^{\mu\nu} &= 0 \\
    \nabla_\mu ^*F^{\mu\nu} &= 0
\end{align}
Where $T^{\mu\nu}$ is the energy-momentum tensor encompassing fluid and electromagnetic contributions, $u^\mu$ is the 4-velocity, and $^*F^{\mu\nu}$ is the dual electromagnetic tensor. 

\textbf{2024-2025 Advancements:} Recent codes integrate \textit{Resistive} Relativistic MHD (RRMHD), allowing for the numerical simulation of magnetic reconnection and current sheets, which accelerate particles to produce superluminal knots.

\section*{3. Particle-in-Cell (PIC) Plasma Simulations}
Used to model the microscopic, collisionless kinetic effects of the plasma. PIC simulations solve the Vlasov-Maxwell system by tracking discrete ``macro-particles''.

\begin{enumerate}
    \item \textbf{Lorentz Force Calculation:} 
    \begin{equation}
        \frac{d(\gamma \mathbf{v})}{dt} = \frac{q}{m} (\mathbf{E} + \mathbf{v} \times \mathbf{B})
    \end{equation}
    \item \textbf{Current/Charge Deposition:} Interpolating particle positions onto a fixed computational grid to calculate charge density $\rho$ and current density $\mathbf{J}$.
    \item \textbf{Maxwell Solvers:} Updating the $\mathbf{E}$ and $\mathbf{B}$ fields on the grid using Finite-Difference Time-Domain (FDTD) Yee algorithms.
\end{enumerate}

\section*{4. Numerical Relativity: Cauchy-Characteristic Evolution (CCE)}
Used to model the permanent spacetime deformations (gravitational wave memory) from massive black hole mergers.

CCE couples inner Cauchy evolution to a ``characteristic'' evolution along outgoing null hypersurfaces extending to future null infinity ($\mathscr{I}^+$). This extracts the non-oscillatory ``memory'' component $\Delta h$ of the gravitational wave strain.

\textbf{2024-2025 Advancements:} The SXS collaboration has developed machine learning surrogate models (NRHybSur3dq8\_CCE) that apply mathematical interpolation to vast databases of CCE simulations to rapidly approximate GW memory signals.

\end{document}
